\subsection{The 144-Component Structure}
\label{subsec:144-structure}

%%%%%%%%%%%%%%%%%%%%%%%%%%%%%%%%%%%%%%%%%%%%%%%%%%%%%%%%%%%%%%%%%%%%%%%%%%%%%%%
% 10.3: The complete 144-component welfare architecture
%%%%%%%%%%%%%%%%%%%%%%%%%%%%%%%%%%%%%%%%%%%%%%%%%%%%%%%%%%%%%%%%%%%%%%%%%%%%%%%

Section~\ref{subsec:three-categories} established three utility categories (INU, KNU, IDN).
Section~\ref{subsec:fepsde-dimensions} established six welfare dimensions (FEPSDE).
This section combines them with time horizons and valence (Gain/Pain) to create
the complete 144-component welfare architecture.

\subsubsection{The Four Structural Elements}

The welfare function is built from four orthogonal elements:

\begin{table}[htbp]
\centering
\caption{The Four Structural Elements}
\label{tab:four-elements}
\renewcommand{\arraystretch}{1.4}
\begin{tabular}{@{}clcl@{}}
\toprule
\textbf{Element} & \textbf{Question} & \textbf{Options} & \textbf{Count} \\
\midrule
Category & For whom? & INU, KNU, IDN & 3 \\
Dimension & In what respect? & F, E, P, S, D, Eco & 6 \\
Time & When? & $t = 0, 1, 2, 3$ & 4 \\
Valence & Gain or Pain? & G, P & 2 \\
\midrule
\multicolumn{3}{r}{\textbf{Total Components:}} & $3 \times 6 \times 4 \times 2 = \mathbf{144}$ \\
\bottomrule
\end{tabular}
\end{table}

%%%%%%%%%%%%%%%%%%%%%%%%%%%%%%%%%%%%%%%%%%%%%%%%%%%%%%%%%%%%%%%%%%%%%%%%%%%%%%%
\subsubsection{Element 1: Utility Categories (3)}
%%%%%%%%%%%%%%%%%%%%%%%%%%%%%%%%%%%%%%%%%%%%%%%%%%%%%%%%%%%%%%%%%%%%%%%%%%%%%%%

From Section~\ref{subsec:three-categories}:

\begin{center}
\renewcommand{\arraystretch}{1.3}
\begin{tabular}{@{}lll@{}}
\toprule
\textbf{Category} & \textbf{Symbol} & \textbf{Question} \\
\midrule
Individual Utility & INU & What do I want for myself? \\
Collective Utility & KNU & What do I want for us? \\
Identity Utility & IDN & Who am I? \\
\bottomrule
\end{tabular}
\end{center}

%%%%%%%%%%%%%%%%%%%%%%%%%%%%%%%%%%%%%%%%%%%%%%%%%%%%%%%%%%%%%%%%%%%%%%%%%%%%%%%
\subsubsection{Element 2: FEPSDE Dimensions (6)}
%%%%%%%%%%%%%%%%%%%%%%%%%%%%%%%%%%%%%%%%%%%%%%%%%%%%%%%%%%%%%%%%%%%%%%%%%%%%%%%

From Section~\ref{subsec:fepsde-dimensions}:

\begin{center}
\renewcommand{\arraystretch}{1.3}
\begin{tabular}{@{}lll@{}}
\toprule
\textbf{Dimension} & \textbf{Symbol} & \textbf{Content} \\
\midrule
Financial & F & Material wealth, economic security \\
Emotional & E & Psychological wellbeing, meaning \\
Physical & P & Health, energy, longevity \\
Social & S & Relationships, belonging, trust \\
Digital & D & Connectivity, access \\
Ecological & Eco & Environmental quality \\
\bottomrule
\end{tabular}
\end{center}

%%%%%%%%%%%%%%%%%%%%%%%%%%%%%%%%%%%%%%%%%%%%%%%%%%%%%%%%%%%%%%%%%%%%%%%%%%%%%%%
\subsubsection{Element 3: Time Horizons (4)}
%%%%%%%%%%%%%%%%%%%%%%%%%%%%%%%%%%%%%%%%%%%%%%%%%%%%%%%%%%%%%%%%%%%%%%%%%%%%%%%

Welfare is not evaluated at a single moment but across time.
We distinguish four horizons:

\begin{definition}[Time Horizons]
\label{def:time-horizons}
\begin{align*}
t = 0 &: \text{\textbf{Instant} --- right now, immediate experience} \\
t = 1 &: \text{\textbf{Short-term} --- days to weeks} \\
t = 2 &: \text{\textbf{Medium-term} --- months to years} \\
t = 3 &: \text{\textbf{Long-term} --- years to lifetime}
\end{align*}
\end{definition}

\paragraph{Why Time Matters.}
Many decisions involve intertemporal trade-offs:

\begin{table}[htbp]
\centering
\caption{Intertemporal Trade-offs}
\label{tab:intertemporal}
\renewcommand{\arraystretch}{1.3}
\begin{tabular}{@{}lll@{}}
\toprule
\textbf{Decision} & \textbf{Short-term} & \textbf{Long-term} \\
\midrule
Education & $U^F_0 \downarrow$ (costs), $U^E_0 \downarrow$ (effort) & $U^F_3 \uparrow$, $U^E_3 \uparrow$ \\
Exercise & $U^E_0 \downarrow$ (discomfort) & $U^P_2 \uparrow$, $U^E_2 \uparrow$ \\
Saving & $U^F_0 \downarrow$ (less consumption) & $U^F_3 \uparrow$ (security) \\
Addiction & $U^E_0 \uparrow$ (pleasure) & $U^{P,E,S,F}_2 \downarrow$ \\
\bottomrule
\end{tabular}
\end{table}

\paragraph{Discounting.}
Future utility is weighted by a discount factor $\delta \in (0, 1)$:
\begin{equation}
\text{Present value of } U_{d,t} = \delta^t \cdot U_{d,t}
\end{equation}

The discount factor may vary by:
\begin{itemize}
\item \textbf{Category:} $\delta^{INU}$ vs. $\delta^{KNU}$ vs. $\delta^{IDN}$
\item \textbf{Dimension:} $\delta_F$ vs. $\delta_P$ (health may be discounted less)
\item \textbf{Context:} $\delta(\Psi_T)$ increases with temporal stability
\end{itemize}

%%%%%%%%%%%%%%%%%%%%%%%%%%%%%%%%%%%%%%%%%%%%%%%%%%%%%%%%%%%%%%%%%%%%%%%%%%%%%%%
\subsubsection{Element 4: Valence --- Gains and Pains (2)}
%%%%%%%%%%%%%%%%%%%%%%%%%%%%%%%%%%%%%%%%%%%%%%%%%%%%%%%%%%%%%%%%%%%%%%%%%%%%%%%

\paragraph{The Prospect Theory Foundation.}
\citet{kahnemantversky1979} demonstrated that people evaluate outcomes 
relative to a reference point, and that losses loom larger than gains:

\begin{itemize}
\item A \$100 loss hurts more than a \$100 gain helps
\item This ``loss aversion'' is approximately 2:1
\item The effect is universal across domains and cultures
\end{itemize}

\paragraph{Gains and Pains Defined.}
For each component, we separate positive deviations (Gains) from 
negative deviations (Pains):

\begin{definition}[Gains and Pains]
\label{def:gains-pains}
Given outcome $x_d$ and reference point $r_d$:
\begin{align}
G_d &= \max(0, x_d - r_d) \quad \text{(Gain: outcome exceeds reference)} \\
P_d &= \max(0, r_d - x_d) \quad \text{(Pain: reference exceeds outcome)}
\end{align}
\end{definition}

Note: For any outcome, either $G_d > 0$ or $P_d > 0$, but not both.

\paragraph{Loss Aversion.}
The utility contribution is asymmetric:
\begin{equation}
U_d = G_d - \lambda_d \cdot P_d
\label{eq:loss-aversion}
\end{equation}

where $\lambda_d > 1$ is the loss aversion coefficient for dimension $d$.

\paragraph{Empirical Estimates of $\lambda$.}

\begin{table}[htbp]
\centering
\caption{Loss Aversion by Dimension}
\label{tab:loss-aversion}
\renewcommand{\arraystretch}{1.3}
\begin{tabular}{@{}lll@{}}
\toprule
\textbf{Dimension} & \textbf{$\lambda$ Estimate} & \textbf{Source} \\
\midrule
Financial ($\lambda_F$) & $\approx 2.0$ & Kahneman \& Tversky (1979) \\
Physical ($\lambda_P$) & $\approx 2.0$--$3.0$ & Health state valuations \\
Social ($\lambda_S$) & $\approx 2.0$--$4.0$ & Rejection sensitivity research \\
Emotional ($\lambda_E$) & $\approx 2.0$ & Affect studies \\
Digital ($\lambda_D$) & $\approx 1.5$--$2.0$ & Emerging evidence \\
Ecological ($\lambda_{Eco}$) & $\approx 1.5$--$2.5$ & Environmental valuation \\
\bottomrule
\end{tabular}
\end{table}

\paragraph{Policy Implication.}
Because $\lambda > 1$, preventing harm is more valuable than creating 
equivalent benefit. A policy that prevents \$100 loss creates approximately
twice the welfare gain of a policy that provides \$100 benefit.

%%%%%%%%%%%%%%%%%%%%%%%%%%%%%%%%%%%%%%%%%%%%%%%%%%%%%%%%%%%%%%%%%%%%%%%%%%%%%%%
\subsubsection{The Complete 144-Component Matrix}
%%%%%%%%%%%%%%%%%%%%%%%%%%%%%%%%%%%%%%%%%%%%%%%%%%%%%%%%%%%%%%%%%%%%%%%%%%%%%%%

\paragraph{Component Notation.}
Each of the 144 components is uniquely identified by four indices:
\begin{equation}
X^i_{d,t} \quad \text{where } X \in \{G, P\}, \; i \in \{INU, KNU, IDN\}, \; d \in \{F,E,P,S,D,Eco\}, \; t \in \{0,1,2,3\}
\end{equation}

\paragraph{Examples.}

\begin{table}[htbp]
\centering
\caption{Examples from the 144-Component Structure}
\label{tab:144-examples}
\renewcommand{\arraystretch}{1.3}
\begin{tabular}{@{}lllll@{}}
\toprule
\textbf{Component} & \textbf{Category} & \textbf{Dim.} & \textbf{Time} & \textbf{Example} \\
\midrule
$G^{INU}_{F,0}$ & Individual & Financial & Instant & Bonus received today \\
$P^{INU}_{P,2}$ & Individual & Physical & Medium & Chronic back pain \\
$G^{KNU}_{S,1}$ & Collective & Social & Short & Team wins award \\
$P^{KNU}_{Eco,3}$ & Collective & Ecological & Long & Climate change damage \\
$G^{IDN}_{E,2}$ & Identity & Emotional & Medium & Pride in achievement \\
$P^{IDN}_{S,0}$ & Identity & Social & Instant & Public humiliation \\
$G^{INU}_{D,1}$ & Individual & Digital & Short & New smartphone \\
$P^{KNU}_{F,2}$ & Collective & Financial & Medium & Company losses \\
$G^{IDN}_{P,3}$ & Identity & Physical & Long & ``I'm still strong at 70'' \\
$P^{INU}_{E,0}$ & Individual & Emotional & Instant & Sudden anxiety \\
$G^{KNU}_{D,2}$ & Collective & Digital & Medium & Community gets broadband \\
$P^{IDN}_{F,1}$ & Identity & Financial & Short & ``I should earn more'' \\
\bottomrule
\end{tabular}
\end{table}

\paragraph{The Matrix Visualization.}

\begin{table}[htbp]
\centering
\caption{The 144-Component Matrix Structure}
\label{tab:144-matrix}
\small
\begin{tabular}{@{}l|cccccc|cccccc@{}}
\toprule
& \multicolumn{6}{c|}{\textbf{GAINS (G)}} & \multicolumn{6}{c}{\textbf{PAINS (P)}} \\
& F & E & P & S & D & Eco & F & E & P & S & D & Eco \\
\midrule
\textbf{INU} & & & & & & & & & & & & \\
\quad $t=0$ & $\bullet$ & $\bullet$ & $\bullet$ & $\bullet$ & $\bullet$ & $\bullet$ & $\bullet$ & $\bullet$ & $\bullet$ & $\bullet$ & $\bullet$ & $\bullet$ \\
\quad $t=1$ & $\bullet$ & $\bullet$ & $\bullet$ & $\bullet$ & $\bullet$ & $\bullet$ & $\bullet$ & $\bullet$ & $\bullet$ & $\bullet$ & $\bullet$ & $\bullet$ \\
\quad $t=2$ & $\bullet$ & $\bullet$ & $\bullet$ & $\bullet$ & $\bullet$ & $\bullet$ & $\bullet$ & $\bullet$ & $\bullet$ & $\bullet$ & $\bullet$ & $\bullet$ \\
\quad $t=3$ & $\bullet$ & $\bullet$ & $\bullet$ & $\bullet$ & $\bullet$ & $\bullet$ & $\bullet$ & $\bullet$ & $\bullet$ & $\bullet$ & $\bullet$ & $\bullet$ \\
\midrule
\textbf{KNU} & & & & & & & & & & & & \\
\quad $t=0$ & $\bullet$ & $\bullet$ & $\bullet$ & $\bullet$ & $\bullet$ & $\bullet$ & $\bullet$ & $\bullet$ & $\bullet$ & $\bullet$ & $\bullet$ & $\bullet$ \\
\quad $t=1$ & $\bullet$ & $\bullet$ & $\bullet$ & $\bullet$ & $\bullet$ & $\bullet$ & $\bullet$ & $\bullet$ & $\bullet$ & $\bullet$ & $\bullet$ & $\bullet$ \\
\quad $t=2$ & $\bullet$ & $\bullet$ & $\bullet$ & $\bullet$ & $\bullet$ & $\bullet$ & $\bullet$ & $\bullet$ & $\bullet$ & $\bullet$ & $\bullet$ & $\bullet$ \\
\quad $t=3$ & $\bullet$ & $\bullet$ & $\bullet$ & $\bullet$ & $\bullet$ & $\bullet$ & $\bullet$ & $\bullet$ & $\bullet$ & $\bullet$ & $\bullet$ & $\bullet$ \\
\midrule
\textbf{IDN} & & & & & & & & & & & & \\
\quad $t=0$ & $\bullet$ & $\bullet$ & $\bullet$ & $\bullet$ & $\bullet$ & $\bullet$ & $\bullet$ & $\bullet$ & $\bullet$ & $\bullet$ & $\bullet$ & $\bullet$ \\
\quad $t=1$ & $\bullet$ & $\bullet$ & $\bullet$ & $\bullet$ & $\bullet$ & $\bullet$ & $\bullet$ & $\bullet$ & $\bullet$ & $\bullet$ & $\bullet$ & $\bullet$ \\
\quad $t=2$ & $\bullet$ & $\bullet$ & $\bullet$ & $\bullet$ & $\bullet$ & $\bullet$ & $\bullet$ & $\bullet$ & $\bullet$ & $\bullet$ & $\bullet$ & $\bullet$ \\
\quad $t=3$ & $\bullet$ & $\bullet$ & $\bullet$ & $\bullet$ & $\bullet$ & $\bullet$ & $\bullet$ & $\bullet$ & $\bullet$ & $\bullet$ & $\bullet$ & $\bullet$ \\
\bottomrule
\end{tabular}
\begin{flushleft}
\small\textit{Note:} Each $\bullet$ represents one component. Total: $3 \times 4 \times 6 \times 2 = 144$ components.
\end{flushleft}
\end{table}

%%%%%%%%%%%%%%%%%%%%%%%%%%%%%%%%%%%%%%%%%%%%%%%%%%%%%%%%%%%%%%%%%%%%%%%%%%%%%%%
\subsubsection{Counting Verification}
%%%%%%%%%%%%%%%%%%%%%%%%%%%%%%%%%%%%%%%%%%%%%%%%%%%%%%%%%%%%%%%%%%%%%%%%%%%%%%%

\paragraph{By Category.}
\begin{align*}
\text{INU components} &= 6 \text{ dimensions} \times 4 \text{ times} \times 2 \text{ valences} = 48 \\
\text{KNU components} &= 6 \times 4 \times 2 = 48 \\
\text{IDN components} &= 6 \times 4 \times 2 = 48 \\
\text{Total} &= 48 + 48 + 48 = 144
\end{align*}

\paragraph{By Dimension.}
\begin{align*}
\text{Financial components} &= 3 \text{ categories} \times 4 \text{ times} \times 2 \text{ valences} = 24 \\
\text{Each other dimension} &= 24 \\
\text{Total} &= 24 \times 6 = 144
\end{align*}

\paragraph{By Time.}
\begin{align*}
\text{Instant } (t=0) &= 3 \times 6 \times 2 = 36 \\
\text{Short-term } (t=1) &= 36 \\
\text{Medium-term } (t=2) &= 36 \\
\text{Long-term } (t=3) &= 36 \\
\text{Total} &= 36 \times 4 = 144
\end{align*}

\paragraph{By Valence.}
\begin{align*}
\text{Gain components} &= 3 \times 6 \times 4 = 72 \\
\text{Pain components} &= 3 \times 6 \times 4 = 72 \\
\text{Total} &= 72 + 72 = 144
\end{align*}

%%%%%%%%%%%%%%%%%%%%%%%%%%%%%%%%%%%%%%%%%%%%%%%%%%%%%%%%%%%%%%%%%%%%%%%%%%%%%%%
\subsubsection{The Symmetry Principle}
%%%%%%%%%%%%%%%%%%%%%%%%%%%%%%%%%%%%%%%%%%%%%%%%%%%%%%%%%%%%%%%%%%%%%%%%%%%%%%%

\paragraph{Structural Identity Across Categories.}
A key feature of the 144-component architecture is that all three 
utility categories have \emph{identical structure}:

\begin{tcolorbox}[colback=blue!5!white, colframe=blue!75!black, title=The Symmetry Principle]
\textbf{INU, KNU, and IDN each have:}
\begin{itemize}
\item 6 FEPSDE dimensions
\item 4 time horizons
\item Gain and Pain components
\item Reference points
\item Loss aversion parameters ($\lambda$)
\item Discount factors ($\delta$)
\end{itemize}

\textbf{The difference lies only in:}
\begin{itemize}
\item How reference points are formed (Section~\ref{subsec:reference-points})
\item Parameter values (which may differ by category)
\end{itemize}

\textbf{This symmetry enables:}
\begin{itemize}
\item Unified calculation logic across all components
\item Clean mathematical formulation
\item Consistent measurement approach
\item Comparable welfare assessment
\end{itemize}
\end{tcolorbox}

%%%%%%%%%%%%%%%%%%%%%%%%%%%%%%%%%%%%%%%%%%%%%%%%%%%%%%%%%%%%%%%%%%%%%%%%%%%%%%%
\subsubsection{Practical Interpretation: A Decision Example}
%%%%%%%%%%%%%%%%%%%%%%%%%%%%%%%%%%%%%%%%%%%%%%%%%%%%%%%%%%%%%%%%%%%%%%%%%%%%%%%

\paragraph{Case: Coal Miner Considering Transition to Coder.}

Consider a coal miner evaluating whether to retrain as a software developer.
We analyze the 144-component structure for this decision:

\begin{table}[htbp]
\centering
\caption{Coal Miner $\rightarrow$ Coder: Selected Components}
\label{tab:coal-miner-example}
\small
\renewcommand{\arraystretch}{1.2}
\begin{tabular}{@{}llllr@{}}
\toprule
\textbf{Component} & \textbf{Cat.} & \textbf{Dim.} & \textbf{Description} & \textbf{Value} \\
\midrule
$G^{INU}_{F,2}$ & INU & Fin. & Higher salary medium-term & $+0.4$ \\
$G^{INU}_{P,2}$ & INU & Phys. & Safer work environment & $+0.2$ \\
$P^{INU}_{S,0}$ & INU & Soc. & Leave mining colleagues now & $-0.2$ \\
\midrule
$P^{KNU}_{S,1}$ & KNU & Soc. & Community loses member & $-0.3$ \\
$G^{KNU}_{F,2}$ & KNU & Fin. & Family better off medium-term & $+0.2$ \\
$G^{KNU}_{Eco,3}$ & KNU & Eco. & Less coal = better for society & $+0.1$ \\
\midrule
$P^{IDN}_{F,0}$ & IDN & Fin. & ``Real men work with hands'' & $-0.4$ \\
$P^{IDN}_{S,0}$ & IDN & Soc. & ``I'm betraying my community'' & $-0.5$ \\
$P^{IDN}_{P,1}$ & IDN & Phys. & ``Miners are tough, coders soft'' & $-0.3$ \\
$P^{IDN}_{E,2}$ & IDN & Emo. & ``Who am I if not a miner?'' & $-0.6$ \\
\bottomrule
\end{tabular}
\end{table}

\paragraph{Aggregation (Preview).}
\begin{align*}
U^{INU} &\approx +0.4 + 0.2 - 0.2 = +0.4 \quad \text{(positive)} \\
U^{KNU} &\approx -0.3 + 0.2 + 0.1 = 0.0 \quad \text{(neutral)} \\
U^{IDN} &\approx -0.4 - 0.5 - 0.3 - 0.6 = -1.8 \quad \text{(strongly negative)}
\end{align*}

Even with positive INU and neutral KNU, the strongly negative IDN 
may dominate the decision---explaining why many miners reject retraining 
despite higher wages.

%%%%%%%%%%%%%%%%%%%%%%%%%%%%%%%%%%%%%%%%%%%%%%%%%%%%%%%%%%%%%%%%%%%%%%%%%%%%%%%
\subsubsection{Summary: The 144-Component Structure}
%%%%%%%%%%%%%%%%%%%%%%%%%%%%%%%%%%%%%%%%%%%%%%%%%%%%%%%%%%%%%%%%%%%%%%%%%%%%%%%

\begin{tcolorbox}[colback=gray!5!white, colframe=gray!75!black, title=Section 10.3 Summary]
\textbf{The 144 components arise from four orthogonal elements:}
\begin{itemize}
\item 3 utility categories (INU, KNU, IDN)
\item 6 FEPSDE dimensions (F, E, P, S, D, Eco)
\item 4 time horizons ($t = 0, 1, 2, 3$)
\item 2 valences (Gain, Pain)
\end{itemize}

$$144 = 3 \times 6 \times 4 \times 2$$

\textbf{Key properties:}
\begin{itemize}
\item \textbf{Symmetric:} All three categories have identical structure
\item \textbf{Comprehensive:} Covers for whom, in what respect, when, and valence
\item \textbf{Prospect-theoretic:} Gains and Pains treated asymmetrically ($\lambda > 1$)
\item \textbf{Temporal:} Explicit time structure with discounting
\end{itemize}

\textbf{Component notation:}
$$X^i_{d,t} \quad \text{where } X \in \{G, P\}, \; i \in \{INU, KNU, IDN\}, \; d \in FEPSDE, \; t \in \{0,1,2,3\}$$

\textbf{Foundation for:}
\begin{itemize}
\item Uniform calculation logic (Section~\ref{subsec:calculation-logic})
\item Reference point analysis (Section~\ref{subsec:reference-points})
\item Aggregation to total welfare (Section~\ref{subsec:aggregation})
\end{itemize}
\end{tcolorbox}

\vspace{1em}
\hrule
\vspace{0.5em}
\noindent\textit{Next section: 10.4 Uniform Calculation Logic}

%%%%%%%%%%%%%%%%%%%%%%%%%%%%%%%%%%%%%%%%%%%%%%%%%%%%%%%%%%%%%%%%%%%%%%%%%%%%%%%
